(1) 非相对论Lorentz磁力大小 \(qvB\) 与速度垂直，提供圆周运动所需向心力 \(mv^2/r\)。由此 \(r=mv/(qB)\)，角频率 \(\omega_c=v/r=qB/m\)，故 \(f_c=qB/(2\pi m)\)，不依赖粒子速度，这正是理想回旋同步的基础。
(2) 每过一次加速间隙，电场做正功 \(qV\)，因此忽略间隙飞行时间及辐射损失时 \(K_N=K_0+NqV\)。再用 \(K_N=mv_N^2/2\) 与 \(r_N=mv_N/(qB)\)，得到 \(r_N=\sqrt{2m(K_0+NqV)}/(qB)\)，半径随能量平方根增长。
(3) 若最后轨道半径不超过 \(R\)，最大非相对论动量为 \(p_{\max}=qBR\)，所以 \(K_{\max}=q^2B^2R^2/(2m)\)。粒子每绕半圈需要加速电压翻转极性，因此电极驱动周期等于完整轨道周期，射频频率为 \(f_{\rm RF}=qB/(2\pi m)\)。
(4) 相对论动量为 \(p=\gamma mv\)，磁场提供横向弯曲使 \(r=p/(qB)\)，角频率变为 \(\omega_{\rm rel}=v/r=qB/(\gamma m)\)。随动能增长 \(\gamma>1\) 增大，所以粒子旋转变慢，而固定射频仍按 \(qB/m\) 驱动，结果逐渐失相，需同步回旋加速器等补偿。
